\section{Controlling a soft residual by a hard residual}
\label{sec:soft_hard_residual}

The residual estimate below concerns two functions of the same positive
matrix. It requires no small global norm of the truncation error. It is
an auxiliary to the soft-truncation argument of Arad, Landau, Vazirani
and Vidick~\cite[lines 1178--1185 and 1195--1255]{Arad2016RigorousRG}.
The scale and Taylor order remain independent.

\begin{lemma}[Norm domination for positive spectral functions]
    \label{lem:positive_cfc_vector_norm_domination}
    \lean{Matrix.norm_cfc_apply_le_of_nonneg_le}
    \leanok
    Let $A$ be a finite Hermitian matrix and let $f,g$ be real functions
    satisfying $0\leq f(x)\leq g(x)$ on its spectrum. Then every vector
    $v$ satisfies
    \begin{align}
        \|f(A)v\|\leq\|g(A)v\|.
    \end{align}
    Neither function is required to be operator monotone.
\end{lemma}
\begin{proof}\leanok
    On the spectrum put $r(x)=f(x)/g(x)$, with $r(x)=0$ when $g(x)=0$.
    The hypotheses give $0\leq r(x)\leq1$ and $f(x)=r(x)g(x)$.
    Functional calculus therefore gives $f(A)=r(A)g(A)$ and
    $\|r(A)\|\leq1$, which proves the assertion. All functions on the
    finite spectrum are continuous.
\end{proof}

\begin{lemma}[Scalar residual comparison at an arbitrary hard cutoff]
    \label{lem:soft_hard_scalar_residual}
    \lean{Real.monotoneOn_softTruncation,
        Real.softTruncation_remainder_le_hard_remainder}
    \leanok
    Let $t>0$ and $m\geq1$. The function $h_{m,t}$ is increasing on
    $[0,\infty)$. For $T,x\geq0$, set
    \begin{align}
        E_{m,t}(T)=\frac{t}{m}\left(\frac{T}{t}\right)^m.
    \end{align}
    Then
    \begin{align}
        0\leq x-h_{m,t}(x)
        \leq x-\min\{x,T\}+E_{m,t}(T).
    \end{align}
\end{lemma}
\begin{proof}\leanok
    Every term of the finite defining series is increasing for $x\geq0$.
    The soft truncation lies below its input. If $x\leq T$, its existing
    polynomial remainder bound is at most $E_{m,t}(T)$. If $x\geq T$,
    monotonicity gives
    \begin{align}
        x-h_{m,t}(x)\leq x-T+T-h_{m,t}(T)
        \leq x-T+E_{m,t}(T).
    \end{align}
    This uses the actual remainder bound below $T$, rather than inferring
    monotonicity of the remainder from monotonicity of the soft function.
\end{proof}

\begin{lemma}[Vector-norm comparison of the residuals]
    \label{lem:soft_hard_vector_residual}
    \lean{Matrix.PosSemidef.norm_softTruncation_residual_apply_le}
    \leanok
    Let $A\geq0$ be a finite positive semidefinite matrix. For every
    $t>0$, $m\geq1$, $T\geq0$ and every vector $v$,
    \begin{align}
        \|(A-h_{m,t}(A))v\|
        \leq\|(A-\min\{A,T\})v\|+E_{m,t}(T)\|v\|.
    \end{align}
    The assertion includes $T=0$ and the zero-dimensional space.
\end{lemma}
\begin{proof}\leanok
    Apply Lemma~\ref{lem:positive_cfc_vector_norm_domination} with
    $f(x)=x-h_{m,t}(x)$ and $g(x)=x-\min\{x,T\}+E_{m,t}(T)$.
    The scalar comparison proves its hypotheses on the nonnegative
    spectrum of $A$. Separate the constant term by functional calculus
    and apply the triangle inequality.
\end{proof}

Thus a hard residual of at most $\rho\|v\|$ on a subspace gives a soft
residual of at most $(\rho+E_{m,t}(T))\|v\|$ there. In particular,
$E_{4,t}(T)=T^4/(4t^3)$, with no substitution $t=m$. The substitution
of actual regional matrices and their physical hard-tail estimates is
a separate step. No reverse soft-to-original comparison, uniform energy
estimate or complete efficient-filter theorem is asserted here.

\begin{lemma}[Residuals restricted to an arbitrary subspace]
    \label{lem:soft_hard_restricted_residual}
    \lean{Matrix.PosSemidef.norm_softTruncation_residual_domRestrict_le}
    \leanok
    Under the hypotheses of Lemma~\ref{lem:soft_hard_vector_residual},
    every subspace $V$ satisfies
    \begin{align}
        \|(A-h_{m,t}(A))|_V\|
        \leq\|(A-\min\{A,T\})|_V\|+E_{m,t}(T).
    \end{align}
\end{lemma}
\begin{proof}\leanok
    For $v\in V$, bound the hard residual by its restricted operator
    norm times $\|v\|$ in Lemma~\ref{lem:soft_hard_vector_residual}.
    Taking the supremum over unit vectors of $V$ gives the assertion.
\end{proof}


\subsection{A quantitative high-energy floor}

\begin{lemma}[Functional calculus and scalar translation]
    \label{lem:cfc_scalar_translation}
    \lean{Matrix.IsHermitian.cfc_sub_smul_one}
    \leanok
    Let $A$ be a finite Hermitian matrix, $e\in\mathbb R$, and
    $g:\mathbb R\to\mathbb R$ any function. Then
    \begin{align}
        g(A-eI)=(x\mapsto g(x-e))(A).
    \end{align}
\end{lemma}
\begin{proof}\leanok
    The functional calculus of $x\mapsto x-e$ is $A-eI$.
    Apply composition in the continuous functional calculus; every
    function on the finite spectrum and its translated image is
    continuous there.
\end{proof}

\begin{lemma}[The soft function above a high-energy threshold]
    \label{lem:soft_truncation_high_energy_floor}
    \lean{Real.softTruncation_high_energy_floor,
        Matrix.IsHermitian.smul_spectralProjectionGE_le_softTruncation}
    \leanok
    Let $t>0$, $m\geq1$ and $R\geq0$. Put
    \begin{align}
        b(R)=R-E_{m,t}(R).
    \end{align}
    Every $x\geq R$ satisfies $h_{m,t}(x)\geq b(R)$.
    Consequently, if $A=A^*\geq eI$, then
    \begin{align}
        b(R)\,\mathbf1_{[e+R,\infty)}(A)
        \leq h_{m,t}(A-eI).
    \end{align}
    The floor $b(R)$ need not be positive for this order comparison.
\end{lemma}
\begin{proof}\leanok
    The polynomial remainder at $R$ gives $h_{m,t}(R)\geq b(R)$.
    Monotonicity gives the same lower bound above $R$.
    On the spectrum of $A$, apply the scalar comparison to $x-e$ when
    $x\geq e+R$, and use $h_{m,t}(x-e)\geq0$ otherwise.
    Functional calculus and Lemma~\ref{lem:cfc_scalar_translation}
    give the projection order.
\end{proof}

\begin{lemma}[A positive denominator for a reverse comparison]
    \label{lem:soft_reverse_barrier_margin}
    \lean{Real.softTruncation_reverse_barrier_margin}
    \leanok
    Let $\eta\geq0$, $m\geq2$, $t\geq2(\eta+8)$,
    $\mu\leq2$ and $c\in\mathbb N$ with $c\leq2$. Put
    \begin{align}
        U=2(\eta+8),\qquad b=U-E_{m,t}(U).
    \end{align}
    Then $U\geq0$, $b\geq\eta+8$, and every $x\geq U$ satisfies
    $h_{m,t}(x)\geq b$. Moreover,
    \begin{align}
        b-(\mu+\eta)-2c\geq2.
    \end{align}
\end{lemma}
\begin{proof}\leanok
    Since $0\leq U/t\leq1$ and $m\geq2$,
    \begin{align}
        E_{m,t}(U)
        \leq\frac{t}{m}\frac Ut
        =\frac Um\leq\frac U2.
    \end{align}
    Hence $b\geq U/2=\eta+8$. The high-energy assertion follows
    from Lemma~\ref{lem:soft_truncation_high_energy_floor}.
    Subtracting $\mu+\eta\leq2+\eta$ and $2c\leq4$ leaves at least
    two.
\end{proof}

The last lemma is a scalar parameter estimate. A physical reverse
comparison must construct the soft Hamiltonian, derive the bound on
its actual ground energy and the cardinality of its actual boundary,
and only then substitute them for $\mu$ and $c$. No reverse residual
or spectral gap is assumed or concluded here. Requiring only a
relative half-threshold error allows fixed Taylor order and scale
linear in the energy width; an absolute unit-error requirement would
be unnecessarily stronger when the width grows.


\subsection{Parameters for arbitrarily small residuals}

\begin{theorem}[A common cutoff and Taylor order]
    \label{thm:soft_residual_parameter_existence}
    \lean{Real.exists_soft_truncation_residual_parameters}
    \leanok
    Let $B\ge0$ and $\varepsilon>0$. There are an integer $m\ge2$,
    a scale $t=3m$ and a cutoff $T\ge B$, with $t\ge B$, such that
    every $b\le B$ satisfies
    \begin{align}
        3\left[32e^{-(T-b)/8}
          +\frac{t}{m}\left(\frac{T}{t}\right)^m\right]
        \le\varepsilon.
        \label{eq:soft_residual_common_parameters}
    \end{align}
    This is a scalar auxiliary for the soft-truncation estimates in
    \cite[lines 1195--1255 and 1373--1420]{Arad2016RigorousRG}.
\end{theorem}
\begin{proof}\leanok
    Put $X=\max\{1,192/\varepsilon\}$ and $T=B+8\log X$.
    Then $32e^{-(T-b)/8}\le32/X\le\varepsilon/6$ whenever $b\le B$.
    Choose an integer $k$ with $2^{-k}<\varepsilon/18$, and then take
    $m\ge\max\{2,k,2T/3\}$. Thus $T/(3m)\le1/2$ and
    \begin{align}
        \frac{t}{m}\left(\frac{T}{t}\right)^m
        =3\left(\frac{T}{3m}\right)^m
        \le3\cdot2^{-m}\le\frac{\varepsilon}{6}.
    \end{align}
    Adding these estimates proves
    (\ref{eq:soft_residual_common_parameters}).
\end{proof}

\begin{corollary}[Both five-region residual budgets]
    \label{cor:soft_cluster_parameter_existence}
    \lean{Real.exists_soft_cluster_parameters}
    \leanok
    Let $g,\tau>0$. There are $m\ge2$, $t=3m$ and a common cutoff $T$
    such that $t,T\ge156+2g$, $T\ge g/40+24$, and
    \begin{align}
        3\left[32e^{-(T-g/40-24)/8}
          +\frac{t}{m}\left(\frac{T}{t}\right)^m\right]
          &\le\frac{\tau g}{10^6},\\
        3\left[32e^{-(T-156-2g)/8}
          +\frac{t}{m}\left(\frac{T}{t}\right)^m\right]
          &\le\frac{\tau g}{10^6}.
    \end{align}
    Both truncation residuals can therefore use the same cutoff.
\end{corollary}
\begin{proof}\leanok
    \uses{thm:soft_residual_parameter_existence}
    Apply the theorem with $B=156+2g$ and
    $\varepsilon=\tau g/10^6$. The two offsets are at most $B$.
\end{proof}

\subsection{A simultaneous count and residual calibration}

The preceding choice of $m$ depends on a prescribed residual tolerance.
It cannot by itself be substituted when that tolerance is later defined
from a count depending on $m$. The next theorem makes all choices from
one natural parameter and verifies the resulting inequalities together.
Its algorithmic scope is discussed in
\cite{gap:alvv16_soft_family_calibration}.

\begin{theorem}[A common natural scale]
    \label{thm:soft_family_calibration}
    \lean{Real.exists_soft_family_calibration}
    \leanok
    Fix $q\in\mathbb N$ with $q\ge2$, $g>0$, and
    $a,b,c\in\mathbb N$. There is $u\in\mathbb N$, $u\ge2$, such that
    the choices
    \begin{align}
        \ell&=u^8,\quad m=u^{10},\quad t=3m,\quad k=u^{16},\quad T=u^9,\\
        D&=(mk+1)(qk)^{36\ell+4\lfloor k/\ell\rfloor},\\
        M&=3t(1+\log m)+4\ell+34,\\
        \tau&=D^{-a}M^{-b},\qquad
        \xi=\min\left\{1,\frac{\tau g^2}{20000(M+g+8)}\right\}
    \end{align}
    satisfy $\ell,k,D\ge1$, $m\ge2$, $M\ge1$, $0<\tau,\xi\le1$,
    $t,T\ge156+2g$, $T\ge g/40+24$, and
    \begin{align}
        3\left[32e^{-(T-g/40-24)/8}
          +\frac tm\left(\frac Tt\right)^m\right]
          &\le\frac{\tau g}{10^6},\\
        3\left[32e^{-(T-156-2g)/8}
          +\frac tm\left(\frac Tt\right)^m\right]
          &\le\frac{\tau g}{10^6},\\
        \left[2e^{-2k\sqrt{(g/2)/(M+\xi+4+g)}}\right]^2
          &\le D^{-12}M^{-c}.
    \end{align}
    The scale $u$ is chosen before any volume or target dimension.
\end{theorem}
\begin{proof}\leanok
    Since $k/\ell=u^8$, direct logarithmic estimates give
    \begin{align}
        \log D&\le u^8(\log2+40\log q+666\log u),\\
        M+R&\le u^{10}(47+R+90\log u)\qquad(R\ge0).
    \end{align}
    The relation $\log u=o(u^p)$ for each $p>0$ now gives, for all
    sufficiently large natural $u$,
    \begin{align}
        M+5+g&\le u^{12},\\
        (a+12)\log D+(b+c)\log M
          +200+g+\left|\log\frac{g}{2000000}\right|
          &\le\frac{u^9}{64}.
    \end{align}
    Increase $u$ further to satisfy the displayed threshold conditions,
    $gu^2\ge1$ and $\log u\ge1$. Exponentiating the second inequality
    bounds the exponential residuals and $9(3u)^{-m}$. The first
    inequality and $gu^2\ge1$ imply
    \begin{align}
        u^9\le4u^{16}\sqrt{\frac{g/2}{M+\xi+4+g}},
    \end{align}
    which proves the squared shrinking bound at the same count.
\end{proof}

This is a new sufficient scalar choice for the soft estimates of
\cite[lines 1262--1268 and 1373--1420]{Arad2016RigorousRG}.
It does not assert the printed fourth-order choice, a quantitative gap
exponent, or a running-time bound.
